% IMO 2026, Problem 3 -- Liu Bang & Xiang Yu stick-cutting game.
% Compile with any standard LaTeX engine (pdflatex/xelatex).
\documentclass[11pt]{article}
\usepackage[margin=1in]{geometry}
\usepackage{amsmath,amssymb,amsthm}
\usepackage{enumitem}
\usepackage{booktabs}

\newtheorem{lemma}{Lemma}
\newtheorem{proposition}{Proposition}
\newtheorem{conjecture}{Conjecture}
\theoremstyle{remark}\newtheorem*{remark}{Remark}

\newcommand{\Sn}{S_{n}}
\newcommand{\Gn}{G_{n}}
\newcommand{\Dstar}{D^{\star}}

\title{IMO 2026, Problem 3 \\[2pt]
\large The Liu Bang--Xiang Yu Stick-Cutting Game}
\author{Solution and partial proof}
\date{}

\begin{document}
\maketitle

\section*{Statement}
Let $n$ be a positive integer. Liu Bang and Xiang Yu have a stick of length $1$.
Liu Bang marks at most $n$ points; then Xiang Yu marks at most $n$ further points
(all marked points distinct). The stick is cut at every marked point, and the
players alternately claim unclaimed pieces, Liu Bang first, each maximizing his
own total length. Determine the largest $c$ that Liu Bang can guarantee.

\section*{Answer}
\[
\boxed{\;c(n)=\dfrac{2^{\,n}}{2^{\,n+1}-1}\;}\qquad
\text{where } \Sn:=2^{n+1}-1.
\]
Equivalently the minimax alternating-sum value is $\Dstar=1/\Sn$, and Liu Bang's
guarantee is $c(n)=\tfrac12(1+\Dstar)=2^n/\Sn$. The first values are
$\tfrac23,\tfrac47,\tfrac{8}{15},\tfrac{16}{31},\tfrac{32}{63}$ (decreasing to
$\tfrac12$).

\paragraph{Proof status.} Proved for $n=1$ (both bounds); \textbf{proved for
$n=2$ (both bounds)}, giving $c(2)=\tfrac47$; lower bound proved for $n=3$
(exhaustive cell enumeration), giving $c(3)\ge\tfrac{8}{15}$; confirmed for
$n\le5$ by exact piecewise-linear search. The general lower bound (Case~B),
the $n=3$ upper bound, and the general upper bound are \emph{conjectured}, with
the obstruction identified; conjectured parts are marked.

\section{Reduction to an alternating-sum game}

After all cuts the pieces satisfy $a_1\ge a_2\ge\cdots\ge a_m>0$, $\sum a_i=1$.

\begin{lemma}[Greedy is optimal]\label{lem:greedy}
In an alternating item-picking game with values $v_1\ge\cdots\ge v_m>0$, greedy play
(always take the largest remaining item) is a subgame-perfect equilibrium for both
players. Player~1 receives $v_1+v_3+v_5+\cdots$; player~2 receives $v_2+v_4+\cdots$.
\end{lemma}
\begin{proof}[Sketch]
Induction on $m$. If player~1 first takes $v_j$ rather than $v_1$, player~2 then
moves first on the remainder, whose odd/even payoff by induction gives
$\text{payoff}(1)-\text{payoff}(j)=(v_1-v_2)+(v_3-v_4)+\cdots\ge0$.
\end{proof}

Hence Liu Bang's total is $\tfrac12(1+D)$ where
\[
D=\sum_{i\ge1}(-1)^{i+1}a_i,\qquad
c(n)=\tfrac12\Bigl(1+\min_{\text{XY}}D\Bigr).
\]

\section{Reformulations}

Work in unnormalized units of total $\Sn$; Liu Bang plays the \emph{geometric
partition} $\Gn=(1,2,4,\ldots,2^n)$, so the target is $\min_{\text{XY}}D=1$.

\paragraph{Even-sum.}
Since $\sum_{\text{odd}}a_i-\sum_{\text{even}}a_i=D$ and
$\sum_{\text{odd}}a_i+\sum_{\text{even}}a_i=\Sn$,
\[
D\ge1\iff \textstyle\sum_{\text{even}}a_i\le \Sn-2^n=2^n-1.
\]

\paragraph{Rank/parity integral.}
With $r(t)=\#\{\text{pieces}\ge t\}$, $D=\int_0^\infty \mathbf 1_{\{r(t)\text{ odd}\}}\,dt$.
A cut splitting a piece of size $s$ into $(m,M)$, $m\le M$, $m+M=s$, changes $r$ by
$+1$ on $(0,m]$, $0$ on $(m,M]$, $-1$ on $(M,s]$; both affected intervals have length
$m$. Thus a cut is a \emph{toggle-pair} of two equal-length intervals.

\section{Lower bound: Liu Bang's geometric construction}

The structural fact is the \emph{$+1$ gap}: $2^k=(2^{k-1}+\cdots+2^0)+1$.

\paragraph{Case A (Xiang Yu does not cut $2^n$).}
Then $2^n$ is the unique largest piece (position~1, odd); every even-position piece
lies in the tail of total $2^n-1$, so $\sum_{\text{even}}\le2^n-1$ and $D\ge1$.

\paragraph{Case B (Xiang Yu cuts $2^n\ge1$ time).}
The fragments of $2^n$ merge with the refined tail $\Gn$ restricted to
$(0,2^{n-1}]$, which is itself the $\Gn_{n-1}$ staircase. The conjectured
\emph{toggle lemma} states that $\le n$ toggle-pairs applied to the geometric
staircase leave at least one unit of odd-parity mass, so $D\ge1$.
Proved for $n\le3$ by exhaustive casework / cell enumeration; conjectured for
$n\ge4$.

\begin{proposition}[$n=1$, proved]\label{prop:n1}
$c(1)=\tfrac23$.
\end{proposition}
\begin{proof}
For two pieces $a\ge b$, $a+b=1$, the geometric play is $(a,b)=(\tfrac23,\tfrac13)$.
Xiang Yu's single cut either splits $a$ (giving $D=\tfrac13$ regardless of where) or
splits $b$ (giving $D\ge\tfrac13$); hence $\min D=\tfrac13$, $c(1)=\tfrac23$. The
matching upper bound $D\le\tfrac13$ is checked by the three regimes
$L\ge\tfrac23$ (halve), $\tfrac12\le L<\tfrac23$ (shave a sliver), and $L=\tfrac12$
(cut one half-piece with $a\le\tfrac16$).
\end{proof}

\begin{proposition}[$n=2$, fully proved]\label{prop:n2}
$c(2)=\tfrac47$.
\end{proposition}
\begin{proof}[Sketch]
\emph{Lower bound} ($c(2)\ge\tfrac47$): exhaustive casework on $k\in\{0,1,2\}$
cuts falling on the piece $4$, with sub-cases for which tail piece the remaining
cut hits; in every region of the (piecewise-linear in the cut positions)
arrangement, the sorted-order merge yields $D\ge1$. Equality $D=1$ is attained at
full halving $4\to2+2,\ 2\to1+1$ (and on whole sub-regions, so the minimizer is
not unique). (Independently confirmed by a $2{,}000{,}000$-sample random sweep:
$0$ violations.)
\emph{Upper bound} ($c(2)\le\tfrac47$): four explicit, universally valid XY
strategies (halve $a_3$, halve $a_1$, match $a_2$, halve both $a_1,a_2$) give
$D\in\{a_1-a_2,\,a_2-a_3,\,|2a_1-1|,\,a_3\}$; a three-line contradiction shows
$\min$ of these $\le1/7$ for every partition $a_1\ge a_2\ge a_3$, $\sum a_i=1$.
Equality at the geometric partition $(\tfrac47,\tfrac27,\tfrac17)$.
\end{proof}

\section{Upper bound: Xiang Yu's strategy}

The matching bound $c(n)\le2^n/\Sn$ says every $\le n+1$-piece partition admits
$\le n$ cuts forcing $D\le1/\Sn$.

\paragraph{$n=1$ (proved).} Treated in Proposition~\ref{prop:n1}.

\paragraph{$n=2$ (proved).} Four explicit XY strategies give $D$ in
$\{a_1-a_2,\;a_2-a_3,\;|2a_1-1|,\;a_3\}$; the key inequality
$\min\le1/7$ (by contradiction: all four $>1/7$ forces $a_3>1/7$,
$a_2>2/7$, $a_1>3/7$, hence $a_1>4/7$, contradicting $a_2+a_3<3/7$)
proves $D\le1/7=1/S_2$ for every three-piece partition.

\paragraph{General $n\ge3$ (conjectured).} The \emph{halving recurrence}: when
$a_1\ge2a_2$, halving $a_1$ gives $D_{\text{new}}=a_1-D_{\text{old}}$. This handles
the \emph{top-heavy} regime $a_1\ge2a_2$ (on $\Gn$ it yields $D=D_0(\Gn_{n-k})$,
reaching $1$ after $n-1$ cuts). The obstruction is the \emph{flat} regime
$a_1<2a_2$: the recurrence breaks, and the global nature of the alternating sum
blocks a naive induction. The conjectured clean route is a
\emph{majorisation/smoothing} principle:
\begin{conjecture}
$f(C):=\min_{\text{XY}}D(C)$ is Schur-maximized uniquely at the geometric
partition; every deformation of a partition toward geometric only increases $f$.
\end{conjecture}
This would close the upper bound from the lower bound. The monotonicity is the
open hard lemma.

\section{Equality}

For $n=1$ and $n=2$ (proved) and conjecturally for all $n$: equality
$c(n)=2^n/\Sn$ is attained in particular when Liu Bang plays geometric and
Xiang Yu plays \emph{full halving}
$2^j\mapsto2^{j-1}+2^{j-1}$ ($j=1,\ldots,n$), producing $2n+1$ pieces
$2^{n-1}(\times2),\ldots,1(\times3)$ with every consecutive pair canceling and
$D=1$ (and on other configurations, e.g.\ the layered-straddle family).

\section*{Summary}

\begin{center}
\begin{tabular}{lcccc}
\toprule
Bound & $n=1$ & $n=2$ & $n=3$ & general $n$\\\midrule
Lower ($c\ge2^n/\Sn$) & proved & proved & proved (cell enum.) & Case~A proved; Case~B conjectured\\
Upper ($c\le2^n/\Sn$) & proved & proved & conjectured & conjectured\\\bottomrule
\end{tabular}
\end{center}

\begin{remark}
The cut budget $\le n$ is essential: with even one extra cut on $\Gn$, Xiang Yu
drives $D$ below $1$. The ``merge inequality'' $D(F\cup T)\ge|F|-|T|=1$ is false for
general tails, so the lower-bound induction must carry the dyadic \emph{structure}
of the refined tail, not merely the numeric bound $D(T)\ge1$.
\end{remark}

\end{document}
